\chapter*{Kata Pengantar}
\addcontentsline{toc}{chapter}{Kata Pengantar}

%=========================================================%
%               TULIS KATA PENGANTAR DI SINI              %
% Jangan Lupa untuk Menghapus Contoh Tulisan di Bawah Ini %
%                                                         %
% [TIP] Sebaiknya jangan dibiasakan menulis gelar atau    %
% jabatan dari nama orang yang dicantumkan di Kata        %
% Pengantar. Penulisan yang tanpa mengurangi rasa hormat  %
% adalah dengan menyebutkan perannya.                     %
%                                                         %
% Contohnya "Gustin Galunggung, selaku profesor yang      %
% mengampu mata kuliah ... dan memberi bimbingan ..."     %
%=========================================================%

Kata pengantar dimaksudkan menyampaikan ungkapan syukur kepada Yang Maha Kuasa dan rasa terima kasih kepada pihak yang berjasa secara langsung dalam penulisan tugas akhir skripsi, serta harapan terkait hasil skripsi Anda.

Puji syukur kehadirat Tuhan Yang Maha Esa atas rahmat dan karunia-Nya, sehingga penulis dapat menyelesaikan skripsi dengan judul ``\judulSebaris'' ini dengan baik dan tepat waktu. Skripsi ini disusun sebagai persyaratan untuk memperoleh gelar Sarjana dalam program studi \prodi.

Penyusunan makalah ini tidak lepas dari bimbingan dan bantuan dari berbagai pihak. Oleh karena itu, pada kesempatan ini penulis ingin mengucapkan terima kasih yang sebesar-besarnya kepada: 1) \namaPembimbingPolos, selaku dosen yang membimbing dan memberi arahan pada proses penyusunan skripsi ini; 2) orang tua dan keluarga yang selalu memberikan dukungan moral dan material; 3) teman-teman komunitas Prodi \prodi\  yang telah bekerja sama,  memberi saran, dan berbagi ilmu; serta 4) semua pihak yang tidak dapat penulis sebutkan satu per satu, atas kontribusi dan dukungannya.

Penulis menyadari bahwa skripsi ini masih jauh dari sempurna. Oleh karena itu, kritik dan saran yang membangun dari pembaca sangat penulis harapkan demi perbaikan di masa mendatang. Semoga skripsi ini dapat memberikan manfaat serta wawasan bagi pembaca.

% --- Tempat, Tanggal, dan Penulis ---
\vfill
\hfill\begin{tabular}{@{}c@{}} 
    \daerahMahasiswa, \tanggalLengkap \vspace{.2cm} \\
    \vspace*{2.1cm} \\ % <-- Setinggi Meterai Rp10.000
    \namaMahasiswa
\end{tabular}
\vfill